\documentclass[../../main.tex]{subfiles}

\begin{document}

\chapter{Pendahuluan dan Orientasi Buku}
\label{chap:pendahuluan}

% =========================================================
% OUTCOME PEMBELAJARAN KHUSUS BAB
% =========================================================

\begin{subcpmkbox}
Bab ini tidak menetapkan Sub-CPMK baru. Bab ini adalah bab
pengikat yang menyiapkan mahasiswa agar dapat mencapai seluruh
Sub-CPMK pada bab-bab berikutnya.

\vspace{0.4em}
\textbf{CPMK yang didukung:} \CPMK{01}, \CPMK{02}, \CPMK{03},
\CPMK{04} \quad\textbullet\quad
\textbf{CPL yang didukung:} \CPL{02}, \CPL{03}
\end{subcpmkbox}

\begin{tujuanbox}
Setelah mempelajari bab ini, mahasiswa mampu:
\begin{enumerate}
    \item menjelaskan tujuan buku ajar dan capaian pembelajaran yang
          hendak dicapai;
    \item menjelaskan keterkaitan buku ajar dengan Rencana
          Pembelajaran Semester (RPS) berbasis \eng{Outcome-Based
          Education} (OBE);
    \item menguraikan jalur capaian pembelajaran dari CPL, CPMK,
          Sub-CPMK, sampai materi dan asesmen;
    \item menggunakan buku ajar secara efektif sesuai alur belajar
          yang disarankan;
    \item menyusun target belajar pribadi dan menyiapkan portofolio
          serta jurnal praktikum.
\end{enumerate}
\end{tujuanbox}

\begin{materibox}
\begin{enumerate}
    \item Tujuan buku ajar dan \eng{learning outcomes} buku.
    \item Keterkaitan buku ajar dengan RPS berbasis OBE.
    \item Jalur capaian pembelajaran: CPL $\rightarrow$ CPMK
          $\rightarrow$ Sub-CPMK $\rightarrow$ materi $\rightarrow$
          asesmen.
    \item Konteks kurikulum OBE dan kedudukan mata kuliah.
    \item Cara menggunakan buku ajar: alur belajar, praktikum, dan
          jurnal refleksi.
\end{enumerate}
\end{materibox}

\subfile{section-01}
\subfile{section-02}
\subfile{section-03}

\section{Contoh: Menelusuri Satu Bab dengan Alur OBE}
\label{sec:contoh-bab1}

Agar alur belajar pada Bagian~\ref{sec:cara-menggunakan} menjadi
konkret, berikut contoh penerapannya pada satu bab, yaitu
Bab~\ref{chap:register} (Register 8086 dan Register Flag) yang
menempati minggu ke-3.

\subsection{Langkah 1: Membaca Sub-CPMK}

Bab tersebut menetapkan \SubCPMK{1.3} dengan indikator: mengidentifikasi
fungsi register umum, register segmen, register \eng{pointer} dan
\eng{index}, serta register \eng{flag} beserta pengaruhnya terhadap
eksekusi program. Dari rumusan ini mahasiswa sudah tahu bahwa yang
akan diuji bukan sekadar kemampuan menghafal nama register, melainkan
kemampuan \emph{memprediksi} perilaku program ketika register berubah.

\subsection{Langkah 2: Membaca Materi dan Menandai Pertanyaan}

Saat membaca, tandai hal yang belum jelas. Contoh pertanyaan yang
wajar muncul pada bab tersebut:

\begin{itemize}
    \item Mengapa \reg{AX} disebut akumulator, dan apakah ia wajib
          dipakai pada operasi aritmatika?
    \item Apa bedanya \reg{SP} dan \reg{BP}? Keduanya menunjuk stack,
          tetapi mengapa perlu dua register?
    \item Jika \flag{CF} dan \flag{OF} sama-sama menunjukkan
          ``kelebihan'', mengapa keduanya dibedakan?
\end{itemize}

Pertanyaan seperti ini dibawa ke kelas. Inilah inti
\eng{flipped classroom}: waktu tatap muka dipakai untuk menjawab
kebingungan, bukan untuk membacakan ulang isi buku.

\subsection{Langkah 3: Menjalankan Contoh Program}

Bagian Contoh Program pada bab tersebut meminta mahasiswa menjalankan
program penjumlahan dan memprediksi nilai flag sebelum program
dijalankan. Tabel~\ref{tab:contoh-prediksi} memperlihatkan bentuk
tabel pengamatan yang diisi setelah program dijalankan.

\begin{table}[H]
    \centering
    \caption{Contoh tabel prediksi dan hasil pengamatan \eng{flag}
    untuk \asm{ADD AL, BL}.}
    \label{tab:contoh-prediksi}
    \small
    \begin{tabular}{@{}cccccc@{}}
        \toprule
        \textbf{AL} & \textbf{BL} & \textbf{Hasil} &
        \textbf{Prediksi ZF} & \textbf{Prediksi CF} &
        \textbf{Keterangan} \\
        \midrule
        \hex{7F} & \hex{01} & \hex{80} & 0 & 0 &
        \textit{overflow} bertanda, \flag{OF} = 1 \\
        \addlinespace
        \hex{FF} & \hex{01} & \hex{00} & 1 & 1 &
        \eng{carry} keluar bit ke-8 \\
        \addlinespace
        \hex{40} & \hex{40} & \hex{80} & 0 & 0 &
        hasil bertanda negatif, \flag{SF} = 1 \\
        \addlinespace
        \hex{10} & \hex{10} & \hex{20} & 0 & 0 &
        hasil biasa, semua flag utama 0 \\
        \bottomrule
    \end{tabular}
\end{table}

\subsection{Langkah 4: Mengisi Jurnal dan Checklist}

Setelah praktikum, mahasiswa mengisi jurnal praktikum sesuai format
Tabel~\ref{tab:format-jurnal}, lalu menandai butir Checklist
Pencapaian Kompetensi pada akhir bab. Butir yang belum bertanda
menjadi daftar tugas belajar sebelum kuis.

\subsection{Langkah 5: Menghubungkan ke Asesmen}

Bagian Evaluasi pada bab tersebut memuat kisi-kisi yang menautkan
setiap soal ke indikator \SubCPMK{1.3}. Dengan membaca kisi-kisi itu
sebelum mengerjakan, mahasiswa mengetahui bobot relatif antara
kemampuan menghafal (\dec{15}\%) dan kemampuan menganalisis
(\dec{85}\%). Informasi ini membantu mengarahkan usaha belajar pada
kemampuan yang benar-benar dinilai.

\begin{infobox}
Perhatikan pola yang berulang: dimulai dari \emph{outcome}, diakhiri
dengan \emph{asesmen}, dan seluruh kegiatan di antaranya adalah
sarana. Pola inilah yang membedakan belajar dalam kerangka OBE dari
belajar menghafal materi.
\end{infobox}

\section{Praktikum Orientasi}
\label{sec:praktikum-bab1}

\begin{praktikumbox}
\textbf{Judul:} Penyiapan portofolio, jurnal praktikum, dan target
belajar pribadi.\\[0.3em]
\textbf{Sub-CPMK:} persiapan untuk 1.1--4.3 \quad\textbullet\quad
\textbf{Alokasi waktu:} 50 menit \quad\textbullet\quad
\textbf{Bentuk:} individu
\end{praktikumbox}

\subsection{Tujuan Praktikum}

Menyiapkan seluruh perangkat administratif dan teknis yang
menunjang pembelajaran sepanjang semester, serta menyusun target
belajar pribadi yang terukur.

\subsection{Alat dan Bahan}

\begin{itemize}
    \item Buku ajar ini beserta dokumen RPS mata kuliah.
    \item Komputer dengan TASM, TLINK, DEBUG, emu8086, DOSBox, dan
          Turbo Debugger.
    \item Direktori kerja, misalnya \file{D:\textbackslash asm}.
\end{itemize}

\subsection{Langkah Kerja}

\begin{enumerate}
    \item \textbf{Menyiapkan struktur direktori kerja.}
    \begin{enumerate}
        \item Buat direktori \file{D:\textbackslash asm}.
        \item Di dalamnya buat subdirektori
              \file{bab03}, \file{bab04}, \ldots{} \file{bab16}.
        \item Setiap program yang Anda tulis disimpan pada
              subdirektori bab yang bersesuaian. Portofolio inilah
              yang dikumpulkan pada akhir semester.
    \end{enumerate}

    \item \textbf{Menyiapkan jurnal praktikum.}
    \begin{enumerate}
        \item Buat berkas \file{jurnal.md} atau buku catatan
              terpisah.
        \item Salin format pada Tabel~\ref{tab:format-jurnal} menjadi
              kerangka jurnal.
        \item Isi satu baris jurnal untuk sesi hari ini sebagai
              latihan.
    \end{enumerate}

    \item \textbf{Memeriksa perangkat praktik.}
    \begin{enumerate}
        \item Periksa ketersediaan TASM, TLINK, dan DEBUG di dalam
              DOSBox.
        \item Periksa ketersediaan emu8086 dan Turbo Debugger.
        \item Baca Bab~VII untuk alur perakitan, pengaitan, dan
              penjalanan program.
    \end{enumerate}

    \item \textbf{Menelusuri buku ajar.}
    \begin{enumerate}
        \item Buka Daftar Isi dan temukan bab yang bersesuaian dengan
              minggu ke-1, ke-8, dan ke-16.
        \item Temukan Lampiran~C (glosarium) dan Lampiran~A (rubrik
              penilaian). Catat halaman masing-masing.
        \item Baca rubrik penilaian praktikum pada Lampiran~A, lalu
              tuliskan satu hal yang akan Anda lakukan agar mencapai
              tingkat tertinggi pada rubrik tersebut.
    \end{enumerate}

    \item \textbf{Menyusun target belajar pribadi.}
    \begin{enumerate}
        \item Tuliskan tiga target belajar untuk semester ini dengan
              rumusan yang terukur. Contoh rumusan yang benar:
              ``Saya dapat menulis program pengurutan assembly
              tanpa melihat contoh pada akhir minggu ke-11.''
        \item Tuliskan satu kebiasaan belajar yang akan Anda jaga,
              misalnya ``menjalankan setiap contoh program dalam
              waktu 24 jam setelah kelas.''
    \end{enumerate}
\end{enumerate}

\subsection{Hasil yang Dikumpulkan}

\begin{enumerate}
    \item Struktur direktori kerja (bukti tangkapan layar).
    \item Berkas jurnal praktikum dengan satu entri terisi.
    \item Lembar target belajar pribadi.
    \item Catatan halaman Daftar Isi, glosarium, dan rubrik penilaian.
\end{enumerate}

\begin{asesmenbox}
\textbf{Asesmen orientasi (bukan komponen nilai, bersifat
prasyarat).} Praktikum ini tidak dinilai secara numerik, tetapi
menjadi \emph{prasyarat} untuk mengikuti praktikum
Bab~\ref{chap:pengantar-assembly}.
Mahasiswa yang belum memiliki struktur direktori dan jurnal tidak
dapat mengumpulkan Tugas 1.

\begin{tabular}{@{}p{7.0cm}p{6.4cm}@{}}
    \toprule
    \textbf{Butir} & \textbf{Kriteria terpenuhi} \\
    \midrule
    Struktur direktori & Subdirektori bab dibuat lengkap dan
    penamaannya konsisten \\
    \addlinespace
    Jurnal praktikum & Format sesuai tabel dan minimal satu entri
    terisi \\
    \addlinespace
    Target belajar & Tiga target terukur dan satu kebiasaan belajar
    dituliskan \\
    \bottomrule
\end{tabular}
\end{asesmenbox}

\section{Latihan}
\label{sec:latihan-bab1}

\begin{enumerate}
    \item \textbf{Orientasi dokumen.} Buka Daftar Isi buku ini.
          Tuliskan judul bab yang membahas masing-masing topik berikut
          beserta nomor babnya: (a) register \eng{flag},
          (b) mode pengalamatan, (c) prosedur dan stack,
          (d) \eng{disassembly}.

    \item \textbf{Pemetaan capaian.} Jelaskan dengan kalimat Anda
          sendiri perbedaan antara CPL, CPMK, dan Sub-CPMK. Berikan
          satu contoh masing-masing yang diambil dari buku ini.

    \item \textbf{Perancangan mundur.} Pada
          Gambar~\ref{fig:backward-design}, materi berada pada urutan
          terakhir. Mengapa demikian? Apa akibatnya bagi materi yang
          tidak mendukung indikator mana pun?

    \item \textbf{Kesepadanan pertemuan.} Anda sedang berada pada
          minggu ke-6 perkuliahan. Bab mana yang seharusnya Anda baca?
          Sub-CPMK apa yang sedang dibangun, dan CPMK apa yang
          didukungnya?

    \item \textbf{Analisis beban.} Perhatikan
          Tabel~\ref{tab:jalur-capaian}. CPMK mana yang mendapat
          bobot terbesar? Menurut Anda, mengapa bobot itu diberikan
          pada CPMK tersebut?

    \item \textbf{Penilaian kinerja.} Sebutkan tiga perbedaan antara
          cara belajar yang cocok untuk asesmen berbasis kinerja dan
          cara belajar yang cocok untuk asesmen berbasis hafalan.
          Kaitkan jawaban Anda dengan komposisi penilaian mata
          kuliah ini.

    \item \textbf{Penyusunan target.} Tuliskan satu target belajar
          untuk diri Anda sendiri dengan syarat: berkaitan dengan
          salah satu Sub-CPMK, dapat diukur, dan memiliki tenggat
          minggu yang jelas.

    \item \textbf{Perencanaan jurnal.} Rancang satu entri jurnal
          praktikum yang lengkap untuk sesi hari ini sesuai
          Tabel~\ref{tab:format-jurnal}. Isi seluruh kolom dengan
          jujur, termasuk kendala yang Anda alami saat pertama kali
          membuka buku ini.
\end{enumerate}

\begin{refleksibox}
Jawab pertanyaan berikut pada jurnal praktikum Anda.
\begin{enumerate}
    \item Apa harapan Anda terhadap mata kuliah ini sebelum membaca
          bab ini? Apakah bab ini mengubah harapan tersebut?
    \item Bagian mana dari kerangka OBE pada
          Gambar~\ref{fig:backward-design} yang paling meyakinkan
          bagi Anda, dan bagian mana yang masih membingungkan?
    \item Kebiasaan belajar apa yang paling perlu Anda ubah agar
          komponen praktikum yang berbobot \dec{70}\% dapat dicapai
          dengan baik?
\end{enumerate}
\end{refleksibox}

\section{Rangkuman}
\label{sec:rangkuman-bab1}

Pada bab ini telah dibahas:

\begin{itemize}
    \item Buku ini bertujuan membuat mahasiswa mampu \emph{bekerja}
          pada level instruksi mesin: menjelaskan cara kerja
          komputer, membaca hubungan instruksi dengan perangkat
          keras, menulis program assembly, dan memperbaiki program
          level rendah.

    \item Empat kemampuan akhir buku dipetakan kepada empat CPMK,
          yang bersama-sama mendukung \CPL{02} dan \CPL{03}.

    \item Pembahasan sengaja dibatasi pada prosesor 8086 mode real
          16-bit. Prosesor ini dipilih karena paling bersih untuk
          mempelajari konsep dasar, dan konsep tersebut berlaku
          langsung pada keluarga x86 modern.

    \item OBE merancang pembelajaran secara mundur: dari CPL menuju
          CPMK, Sub-CPMK, asesmen, baru kemudian materi. Materi
          hanya dianggap perlu bila mendukung indikator yang dapat
          diukur.

    \item Terdapat \dec{14} Sub-CPMK yang dipetakan ke
          \dec{16} minggu perkuliahan. Setiap bab materi buku ini
          bersesuaian dengan satu atau dua minggu perkuliahan.

    \item Setiap bab materi memiliki kerangka baku: Sub-CPMK, tujuan
          pembelajaran, materi pokok, contoh program, praktikum,
          latihan, rangkuman, dan evaluasi.

    \item Penilaian mata kuliah berbobot \dec{70}\% praktik dan
          \dec{30}\% teori, sehingga portofolio program dan jurnal
          praktikum menjadi bukti utama pencapaian kompetensi.

    \item Tiga perangkat lunak yang dipakai sepanjang semester adalah
          Emu8086, GUI Turbo Assembler, dan DOSBox.
\end{itemize}

\begin{checklistbox}
Tandai setiap butir yang sudah Anda kuasai.

\begin{itemize}
    \item[\checkbox] Saya dapat menjelaskan tujuan buku ini dan empat
          kemampuan akhir yang hendak dicapai.
    \item[\checkbox] Saya dapat menjelaskan jalur CPL $\rightarrow$
          CPMK $\rightarrow$ Sub-CPMK $\rightarrow$ materi
          $\rightarrow$ asesmen.
    \item[\checkbox] Saya dapat menemukan bab yang bersesuaian dengan
          minggu perkuliahan tertentu.
    \item[\checkbox] Saya sudah memiliki struktur direktori kerja
          untuk menyimpan portofolio program.
    \item[\checkbox] Saya sudah memiliki jurnal praktikum dengan
          format yang benar.
    \item[\checkbox] Saya sudah menyusun target belajar pribadi yang
          terukur.
    \item[\checkbox] Saya mengetahui letak glosarium dan rubrik
          penilaian pada lampiran.
\end{itemize}

\vspace{0.4em}
Seluruh butir di atas bersifat prasyarat. Pastikan semuanya bertanda
sebelum memasuki Bab~\ref{chap:landasan-teori}.
\end{checklistbox}

\begin{infobox}
\textbf{Rujukan bab.} Kerangka OBE pada bab ini mengikuti struktur
silabus dan RPS mata kuliah. Untuk memperluas wawasan mengenai
organisasi komputer yang menjadi latar mata kuliah ini, lihat
\cite{stallings2016} dan \cite{tanenbaum2013}. Untuk gambaran
praktis mengenai cara kerja prosesor 8086, lihat \cite{intel1979}.
\end{infobox}

\section{Evaluasi}
\label{sec:evaluasi-bab1}

Evaluasi pada bab ini berbentuk \textbf{kuis diagnostik}. Kuis ini
tidak dinilai, tetapi wajib dikerjakan karena hasilnya menjadi
titik tolak bagi dosen dan mahasiswa dalam merencanakan pembelajaran
sepanjang semester.

\subsection{Soal Pilihan Ganda}

\begin{enumerate}
    \item Pada perancangan OBE, urutan yang benar adalah \dots
          \begin{enumerate}
              \item materi $\rightarrow$ asesmen $\rightarrow$ CPMK
                    $\rightarrow$ CPL
              \item CPL $\rightarrow$ CPMK $\rightarrow$ Sub-CPMK
                    $\rightarrow$ asesmen $\rightarrow$ materi
              \item CPMK $\rightarrow$ materi $\rightarrow$ CPL
                    $\rightarrow$ asesmen
              \item asesmen $\rightarrow$ materi $\rightarrow$ CPL
                    $\rightarrow$ CPMK
          \end{enumerate}

    \item Indikator yang lebih terukur dan dapat diuji, yang
          diturunkan dari CPMK, disebut \dots
          \begin{enumerate}
              \item CPL \item CPMK
              \item Sub-CPMK \item rubrik
          \end{enumerate}

    \item Bab yang membahas segmentasi memori dan alamat fisik adalah
          \dots
          \begin{enumerate}
              \item Bab~V \item Bab~VI
              \item Bab~VII \item Bab~VIII
          \end{enumerate}

    \item CPMK yang didukung oleh Bab~X sampai Bab~XIII adalah \dots
          \begin{enumerate}
              \item \CPMK{01} \item \CPMK{02}
              \item \CPMK{03} \item \CPMK{04}
          \end{enumerate}

    \item Komposisi penilaian mata kuliah ini adalah \dots
          \begin{enumerate}
              \item \dec{50}\% teori dan \dec{50}\% praktik
              \item \dec{30}\% teori dan \dec{70}\% praktik
              \item \dec{70}\% teori dan \dec{30}\% praktik
              \item \dec{100}\% ujian tertulis
          \end{enumerate}
\end{enumerate}

\subsection{Soal Uraian}

\begin{enumerate}
    \item Jelaskan mengapa pada OBE materi berada pada urutan terakhir
          dalam perancangan. Sertakan satu contoh materi yang mungkin
          \emph{tersisih} karena tidak mendukung indikator apa pun.

    \item Anda adalah mahasiswa yang baru masuk pada minggu ke-9.
          Bab mana yang harus Anda baca, dan Sub-CPMK apa yang sedang
          dibangun? Jelaskan alasan jawaban Anda.

    \item Bandingkan cara Anda mempersiapkan diri untuk komponen
          UTS (berbobot \dec{20}\%) dan komponen praktikum (berbobot
          \dec{20}\%). Apa perbedaan strategi belajarnya?

    \item Tuliskan rencana belajar Anda untuk empat minggu pertama
          dalam bentuk tabel: minggu, bab yang dibaca, Sub-CPMK yang
          ditargetkan, dan bukti belajar yang akan dikumpulkan.
\end{enumerate}

\subsection{Kisi-kisi Kuis Diagnostik}

\begin{table}[H]
    \centering
    \caption{Kisi-kisi kuis diagnostik Bab~\ref{chap:pendahuluan}.}
    \label{tab:kisi-bab1}
    \small
    \begin{tabular}{@{}p{5.6cm}p{5.0cm}c@{}}
        \toprule
        \textbf{Aspek yang diperiksa} & \textbf{Bentuk soal} &
        \textbf{Bobot} \\
        \midrule
        Pemahaman kerangka OBE & Pilihan ganda 1, 2; Uraian 1 &
        30\% \\
        \addlinespace
        Kemampuan menavigasi buku dan RPS & Pilihan ganda 3, 4;
        Uraian 2 & 30\% \\
        \addlinespace
        Pemahaman sistem penilaian & Pilihan ganda 5; Uraian 3 &
        20\% \\
        \addlinespace
        Kemampuan merencanakan belajar & Uraian 4 & 20\% \\
        \bottomrule
    \end{tabular}
\end{table}

\begin{checklistbox}
\textbf{Checklist Kesiapan Mengikuti Mata Kuliah}

Seorang mahasiswa dinyatakan siap mengikuti mata kuliah ini apabila:
\begin{itemize}
    \item[\checked] dapat menjelaskan tujuan buku dan keempat
          kemampuan akhirnya;
    \item[\checked] dapat menjelaskan jalur capaian pembelajaran dari
          CPL sampai asesmen;
    \item[\checked] dapat menemukan bab yang bersesuaian dengan
          minggu perkuliahan;
    \item[\checked] memiliki direktori kerja, jurnal praktikum, dan
          target belajar pribadi;
    \item[\checked] mengetahui komposisi penilaian dan letak rubrik
          penilaiannya.
\end{itemize}

\vspace{0.4em}
Bila ada butir yang belum bertanda, ulangi Bagian~\ref{sec:cara-menggunakan}
sebelum melanjutkan.
\end{checklistbox}


\end{document}
